\par\Needspace{10\baselineskip}\pdfbookmark[1]{Matthiessen. 1819 Januar 30.}{matthiessen1819}\medskip\hrule\smallskip
\begin{center}
Göttingische gelehrte Anzeigen. 1819 Januar 30.
\end{center}
\smallskip\hrule\medskip

[E. A. \textsc{Matthiessen}.] \textit{Tafel zur bequemern Berechnung des Logarithmen der
Summe oder Differenz zweier Grössen, welche selbst nur durch ihre Logarithmen gegeben sind.} Altona 1818. Bei J. F. Hammerich. Einleitung 33 S. Die Tafeln
212 S. in Quart. (Titel und Einleitung auch in lateinischer Sprache.)

Bei etwas umfangreicheren Rechnungen ist jede Abkürzung wertvoll; auch Hilfsmittel, die bei einmaliger Anwendung nur einen kleinen Vorteil bieten, werden durch häufig wiederkehrenden Gebrauch wichtig. Ein solches Hilfsmittel ist eine 1812 erstmals in der Monatlichen Correspondenz veröffentlichte Tafel, mit der sich aus den Logarithmen zweier Größen unmittelbar die Logarithmen ihrer Summe oder Differenz ableiten lassen: Man erreicht mit einmaligem Nachschlagen, wozu sonst dreimaliges oder wenigstens zweimaliges nötig wäre. Mit einer solchen Tafel wird man manche \SourceLeaf{260} Formeln mit Vorteil in ihrer ursprünglichen Gestalt beibehalten können, denen man sonst durch Einführung von Hilfswinkeln eine für die Rechnung bequemere Form zu geben sucht. Die erwähnte Tafel war nur für Rechnungen mit fünf Dezimalstellen bestimmt: Gerade weil solche Rechnungen am häufigsten vorkommen, war dies das nächstliegende Bedürfnis. Der bei ihrer Veröffentlichung geäußerte Wunsch, jemand möge sich die Mühe machen, eine ähnliche Tafel in größerem Umfang und mit sieben Dezimalstellen zu berechnen, hat die vorliegenden Tafeln veranlasst. Ihr Verfasser, Herr \textsc{Matthiessen}, hat sich mit dieser mühsamen Arbeit ein Recht auf den Dank aller erworben, die viel mit logarithmischen Rechnungen zu tun haben. Die Vorrede berichtet über die vom Verfasser zur Berechnung verwendete Methode und über seine lobenswerte Sorgfalt, die Tafel durchgehend auch in der letzten Ziffer zuverlässig zu machen. Bei Tafeln dieser Art ist auch die äußere Einrichtung keineswegs gleichgültig. Der Verfasser hatte zunächst eine Anordnung im Sinn, die die Tafel auf den kleinstmöglichen Raum zusammendrängen sollte. Da die Tafel aber ausschließlich dazu bestimmt ist, die Rechnungen zu \textit{erleichtern}, würde dieser Zweck ganz verfehlt, wenn ihre Anordnung zu verwickelt wäre und mühsame Aufmerksamkeit erforderte. Nach reiflicher Überlegung entschloss sich der Verfasser daher mit Recht, jene zwar sinnreiche Anordnung aufzugeben und stattdessen eine einfachere zu wählen, obwohl sich der Umfang des Bandes dadurch erheblich vergrößerte. Vielleicht wird mancher Benutzer mit uns wünschen, der Verfasser wäre darin noch etwas weiter gegangen und hätte die Spalten \(B\) und \(C\) jeweils vollständig abdrucken lassen; ihre letzten vier, beiden gemeinsamen Ziffern stehen nur einmal da. Das hätte die Übersicht noch bequemer gemacht, ohne das Format zu vergrößern, wenn eine etwas kleinere Schrift gewählt worden wäre. Das wäre ohne Nachteil, vielleicht sogar mit Vorteil für das gefällige Aussehen möglich gewesen. Auch die fast zu strenge Sparsamkeit bei den Ziffern, von denen in der Regel nur die vierte abgedruckt ist, solange die vorangehenden unverändert bleiben, und bei den Proportionalteilen, die immer nur einmal angegeben sind und daher manchmal ein Zurück- oder Vorblättern erfordern, beeinträchtigt die Bequemlichkeit etwas. Schließlich hätten wir gewünscht, der letzte Teil der Tafel auf S. 212 wäre in einem zehnmal so großen Umfang gegeben, um die zweiten Differenzen unmerklich zu machen; dazu wäre \SourceLeaf{261} nur eine zusätzliche Seite nötig gewesen. Alle diese Bemerkungen, die zum Teil auch auf individueller Gewöhnung beruhen mögen, sollen das Verdienst dieser Arbeit keineswegs schmälern. Es wird gewiss von allen anerkannt, die Gelegenheit haben werden, sie zu benutzen.

